\section{The End: Death at Rome}

\subsection{What the first witnesses say}

No New Testament text narrates Paul's death. The earliest statement of it
comes from the church of Rome itself, in the letter to Corinth known as
1~Clement, conventionally dated c.~95--96---one generation after the
event, from the city where it happened. Among the ``most recent spiritual
heroes'' persecuted ``through envy'':

\begin{studyframe}\small
``Owing to envy, Paul also obtained the reward of patient endurance, after
being seven times thrown into captivity, compelled to flee, and stoned.
After preaching both in the east and west, he gained the illustrious
reputation due to his faith, having taught righteousness to the whole
world, and come to the extreme limit of the west, and suffered martyrdom
under the prefects. Thus was he removed from the world, and went into the
holy place, having proved himself a striking example of patience''
(1~Clem.\ 5).
\end{studyframe}

Every clause matters. ``Seven times thrown into captivity'' exceeds
anything Acts narrates---independent memory, or rhetoric, of a much-arrested man.
``Suffered martyrdom under the prefects'' places the death under Roman
authority; paired, in the same breath, with Peter and with ``a great
multitude of the elect'' who suffered with them, the natural referent is
Nero's Roman persecution after the fire of 64, when, as Tacitus records,
``an immense multitude was convicted'' and mocked, burned, and crucified
(\work{Annals} 15.44). Within another generation the pairing is assumed:
Ignatius of Antioch, en route to his own death, tells the Romans, ``I do
not, as Peter and Paul, issue commandments unto you'' (\work{Romans}~4),
and tells the Ephesians they are ``initiated into the mysteries of the
Gospel with Paul, the holy, the martyred, the deservedly most happy''
(\work{Ephesians}~12). Polycarp holds up ``the blessed and glorified
Paul,'' with Ignatius and others, among those now ``in their due place in
the presence of the Lord, with whom also they suffered''
(\work{Philippians}~3; 9). By c.~170 Dionysius of Corinth says Peter and
Paul ``taught together in like manner in Italy, and suffered martyrdom at
the same time''; c.~200 the Roman presbyter Gaius points to the monuments
themselves---``if you will go to the Vatican or to the Ostian way, you
will find the trophies of those who laid the foundations of this
church''---and Eusebius, who preserves both, states the received sum:
``It is, therefore, recorded that Paul was beheaded in Rome itself, and
that Peter likewise was crucified under Nero'' (\work{Church History}
2.25.5--8; cf.\ 3.1.1, from Origen). The chain is unbroken, geographically
Roman, and unrivaled: no other city ever claimed Paul's death or grave.
Martyrdom at Rome under Nero is as close to certain as an unnarrated
first-century event can be. The mode---the sword---is stated by
Tertullian c.~200 (Rome, ``where Paul wins his crown in a death like
John's,'' i.e.\ the Baptist's, \work{Prescription}~36) and by the roughly
contemporary \work{Acts of Paul}; it fits the citizen of Acts, and no
ancient source contradicts it.

\subsection{``The extreme limit of the west'' and the Spanish question}

Did Paul reach Spain? He intended to: ``When I shall begin to take my
journey into Spain, I hope that as I pass, I shall see you\ldots\ I will
come by you into Spain'' (Rom 15:24, 28). Two early texts can be read as
saying he did. 1~Clement's ``extreme limit of the west'' would most
naturally mean Spain to a Roman writer---unless it is generous rhetoric
for Rome itself, written as it is in Rome, in a sentence built of
superlatives; the phrase is genuinely ambiguous, and Benedict~XVI's
catechesis states the case exactly: ``Whether this is a reference to a
voyage in Spain undertaken by Paul is open to discussion. There is no
certainty on it'' (General Audience, 4 February 2009). The Muratorian
fragment, listing what Luke omitted, names ``the passion of Peter, and
also of the journey of Paul, when he went from the city---Rome---to
Spain'': a Roman document of (conventionally) c.~170--200 that assumes the
journey---or assumes Romans 15 fulfilled. Later tradition kept the trip
alive because it needed it: Eusebius reports---as report (``it is
said'')---that after the two years of Acts ``the apostle was sent again
upon the ministry of preaching, and that upon coming to the same city a
second time he suffered martyrdom'' (\work{Church History} 2.22), reading
2~Timothy's ``At my first answer no man stood with me\ldots\ and I was
delivered out of the mouth of the lion'' (2~Tim 4:16--17) as the record of
a first, survived trial; Jerome hardens report into fact: ``Paul was
dismissed by Nero, that the gospel of Christ might be preached also in the
West'' (\work{De viris illustribus}~5). The sober balance: intention
certain; execution unproven; a release c.~62, a western (and/or renewed
eastern, so the Pastorals' itineraries) ministry, and a second Roman
arrest form a coherent, early-attested, and unverifiable reconstruction.
The rival reading---one continuous imprisonment from Acts 28 to death,
with 1~Clement's ``limit of the west'' as Rome---is equally coherent.
This study asserts neither.

\subsection{The year and the day}

The persecution gives the floor date of 64. The ceiling is Nero's death in
June 68. Within that window the ancient chronographic tradition placed
Paul late: Jerome, ``in the fourteenth year of Nero on the same day with
Peter, he was beheaded at Rome for Christ's sake and was buried in the
Ostian way'' (\work{De viris illustribus}~5)---i.e.\ 67/68; the 1911
Catholic Encyclopedia collects the variants (Epiphanius Nero's twelfth
year, Euthalius the thirteenth, Jerome the fourteenth). ``On the same day
with Peter'' is calendar piety more than chronology: the oldest witness to
the pairing, Dionysius, says only ``at the same time'' (Greek \textit{kata
ton auton kairon}, which need mean no more than ``about the same time''),
and 29 June, the joint feast, is first attested as a Roman cult date in
the mid-third century, not as anyone's death date. ``Between 64 and 68,
probably toward the later years if the two-imprisonment tradition has
substance,'' is all the evidence yields; Benedict~XVI's martyrdom
catechesis states the same range on the same grounds. The scene of the
beheading---at the Aquae Salviae, south of the Ostian road---enters the
record only in late martyrdom romances, and belongs with the legends
below.
